Hustota odborových organizací v~ČR prošla od roku 1993 dramatickým poklesem: z~přibližně \SI{70}{\percent} v~době transformace na současných asi \SI{11}{\percent}, přičemž trend není zastaven a~ČR se přibližuje hodnotám typickým pro Francii (\SI{8}{\percent}), kde je však nízká hustota kompenzována zákonným rozšiřováním závaznosti \ac{KS} na celé odvětví.
Tento propad není jen statistickou kuriozitou --- představuje fundamentální ztrátu vyjednávací síly zaměstnanců: s~každým procentním bodem hustoty klesá schopnost odborů prosazovat nadstandardní mzdová ujednání i~jiné nesení podmínek, jdoucí nad zákonné minimum.
AT a~DK udržují hustotu na \SIrange{55}{65}{\percent}, přičemž právě tyto země vykazují nejnižší mzdovou disperzí a~nejvyšší podíl pracovníků pokrytých kolektivní smlouvou --- doklady pro kauzalitu ve smyslu Freemana a~Medoffa (1984).
Obnova sociálního dialogu v~ČR musí tedy řešit nejen právní rámec, ale především důvěru zaměstnanců v~odbory jako instituci, která je schopna přinést hmatatelné mzdové a~pracovní benefity.

Hustota odborových organizací v~ČR prošla od roku 1993 dramatickým poklesem: z~přibližně \SI{70}{\percent} v~době transformace na současných asi \SI{11}{\percent}, přičemž trend není zastaven a~ČR se přibližuje hodnotám typickým pro Francii (\SI{8}{\percent}), kde je však nízká hustota kompenzována zákonným rozšiřováním závaznosti \ac{KS} na celé odvětví.
Tento propad není jen statistickou kuriozitou --- představuje fundamentální ztrátu vyjednávací síly zaměstnanců: s~každým procentním bodem hustoty klesá schopnost odborů prosazovat nadstandardní mzdová ujednání i~jiné nesení podmínek, jdoucí nad zákonné minimum.
AT a~DK udržují hustotu na \SIrange{55}{65}{\percent}, přičemž právě tyto země vykazují nejnižší mzdovou disperzí a~nejvyšší podíl pracovníků pokrytých kolektivní smlouvou --- doklady pro kauzalitu ve smyslu Freemana a~Medoffa (1984).
Obnova sociálního dialogu v~ČR musí tedy řešit nejen právní rámec, ale především důvěru zaměstnanců v~odbory jako instituci, která je schopna přinést hmatatelné mzdové a~pracovní benefity.

Hustota odborových organizací v~ČR prošla od roku 1993 dramatickým poklesem: z~přibližně \SI{70}{\percent} v~době transformace na současných asi \SI{11}{\percent}, přičemž trend není zastaven a~ČR se přibližuje hodnotám typickým pro Francii (\SI{8}{\percent}), kde je však nízká hustota kompenzována zákonným rozšiřováním závaznosti \ac{KS} na celé odvětví.
Tento propad není jen statistickou kuriozitou --- představuje fundamentální ztrátu vyjednávací síly zaměstnanců: s~každým procentním bodem hustoty klesá schopnost odborů prosazovat nadstandardní mzdová ujednání i~jiné nesení podmínek, jdoucí nad zákonné minimum.
AT a~DK udržují hustotu na \SIrange{55}{65}{\percent}, přičemž právě tyto země vykazují nejnižší mzdovou disperzí a~nejvyšší podíl pracovníků pokrytých kolektivní smlouvou --- doklady pro kauzalitu ve smyslu Freemana a~Medoffa (1984).
Obnova sociálního dialogu v~ČR musí tedy řešit nejen právní rámec, ale především důvěru zaměstnanců v~odbory jako instituci, která je schopna přinést hmatatelné mzdové a~pracovní benefity.

Hustota odborových organizací v~ČR prošla od roku 1993 dramatickým poklesem: z~přibližně \SI{70}{\percent} v~době transformace na současných asi \SI{11}{\percent}, přičemž trend není zastaven a~ČR se přibližuje hodnotám typickým pro Francii (\SI{8}{\percent}), kde je však nízká hustota kompenzována zákonným rozšiřováním závaznosti \ac{KS} na celé odvětví.
Tento propad není jen statistickou kuriozitou --- představuje fundamentální ztrátu vyjednávací síly zaměstnanců: s~každým procentním bodem hustoty klesá schopnost odborů prosazovat nadstandardní mzdová ujednání i~jiné nesení podmínek, jdoucí nad zákonné minimum.
AT a~DK udržují hustotu na \SIrange{55}{65}{\percent}, přičemž právě tyto země vykazují nejnižší mzdovou disperzí a~nejvyšší podíl pracovníků pokrytých kolektivní smlouvou --- doklady pro kauzalitu ve smyslu Freemana a~Medoffa (1984).
Obnova sociálního dialogu v~ČR musí tedy řešit nejen právní rámec, ale především důvěru zaměstnanců v~odbory jako instituci, která je schopna přinést hmatatelné mzdové a~pracovní benefity.

\input{texparts/python/union_density_trend}



